\chapter{Discussion}\label{ch:discussion}

Chapter~\ref{ch:experiments} ends with a problem. On the same
Russo-Ukrainian conflict, four measurement instruments do not agree
on which language the cross-lingual bias points toward: anchor
cosine pulls English-ward at $+0.175$ and paired-edit stance
projection compresses to $+0.033$, both on the real-event corpus;
the prose-stance judge tilts Ukrainian-ward at $-0.084$ on the
imaginary corpus; and
\citet{li2024borderlines}'s discrete MCQ flips Russian-ward in
vanilla GPT-4, an external published result on a model outside the
matrix. A multilingual-bias finding that depends on which
instrument was used is either an instrument story or a phenomenon
story, and the chapter has to decide. This discussion argues for the
second reading: the four instruments measure different things on the
same model, and the disagreement between them is substantive rather
than an artefact of noise.

\section{Robust findings}\label{sec:robust_findings}

The headline cross-lingual ingroup-pull finding is more thoroughly
stress-tested than the original thesis proposal required. The
\emph{sign} of the pull --- responses embed closer to the own-language
Wikipedia anchor than to other-language anchors --- holds on the
English diagonal for all fourteen cosine-eligible providers and on
all three diagonals for thirteen of the fourteen, across six
regulatory regimes
(\S\ref{sec:exp_014}); at the per-language level it holds across
all four embedders spanning four
ecosystems (\S\ref{sec:exp_015}), under a 2$\times$ change in token
budget (\S\ref{sec:exp_002}), under unsupervised clustering that
recovers the design with no axis supervision (\S\ref{sec:exp_016}),
across five conflict families spanning six writing systems
(with the Taiwan-strait Chinese cell as the single marginal
exception, \S\ref{sec:exp_006}), and across eight bilateral
disputes and 31
events spanning a millennium, against two settled-dispute baselines
that sit at the floor of the contested range and remain below it
under the length-matched anchor control --- narrowly at the 2022
cell, whose matched margin is $+0.007$
(\S\ref{sec:exp_007}). On
the Falklands the magnitude collapses to $+0.05$,
functioning as a falsification anchor for the pipeline rather than a
counterexample to it. The sign is the load-bearing claim. The
settled baselines calibrate the raw floor, with the qualification
\S\ref{sec:exp_007} develops: the settled pairs' mutual
intelligibility mechanically compresses their attainable pull, so
the raw settled-vs-contested magnitude separation is partly an
anchor-geometry effect and cannot by itself attribute the contested
families' higher raw pull to contestation.

One framing carried through earlier chapters has to be given up
here. The Falklands near-null was read as an instance of
\citet{oeberst2020wikipedia}'s recency--strength gradient, and more
broadly the expectation was that a model reproducing its training
distribution would inherit that gradient along with the bias. The
event ladders test this directly and it does not hold. Within
conflict families the correlation between event date and pull splits
three positive against three negative, and events from 1453 and 1772
return pull well above the 1982 Falklands figure. The model
inherits the \emph{direction} of Wikipedia's ingroup bias without
inheriting its \emph{temporal structure} --- which is a more
interesting result than the confirmation would have been, because it
separates the source-side mechanism of \S\ref{sec:mechanisms} from
the simple corpus-reproduction story. What produces the Falklands
near-null remains open; recency is no longer an available answer.

The finding that survives is thus an ingroup pull that is
\emph{topical} rather than \emph{positional}, and this is what
reconciles it with \citet{durmus2023globalopinionqa}'s Linguistic
Prompting null rather than contradicting it. Their result is that
prompting in a language does not move a model toward that
language community's opinions; the result here is that prompting in
a language does move the response toward that language's
encyclopedic treatment of the event. Both can hold at once because
resembling a corpus and adopting its side are different things ---
the argument developed in \S\ref{sec:four_instruments}. The
contribution is to have generalised the comparison from opinion
elicitation to free-form generation about contested events, scaled it
from one provider to a fifteen-provider matrix, and grounded it
against both a falsification anchor and two settled-dispute
baselines that the original study did not include.

\section{Falsified framings}\label{sec:falsified_framings}

Three stronger framings the project carried at various points do not
survive.

The \emph{magnitude ordering} EN${\gg}$RU${>}$UK is OpenAI-specific.
Under three other embedders the pattern flattens to
EN${\approx}$UK${>}$RU, and the Russian-trained Yandex embedder gives
Ukrainian responses \emph{more} pull than Russian responses
($+0.090$ vs $+0.042$). The ordering is a property
of the OpenAI embedding space, not of the responses, and the thesis
prose treats it as such throughout.

The literal ``models default to a Russian framing on imaginary
contested events'' reading is falsified on three independent
measurement arms --- cosine divergence, prose register, and the
stance judge --- none of which returns a Russian-ward mean. What is
there instead is the two-mode claim of \S\ref{sec:two_modes}.

The lexical-versus-content cut from
\S\ref{sec:exp_001_debiased} --- ``UK pull is lexical, EN pull is
content-driven'' --- is sharper than the data support. Under Alibaba
and Yandex the Ukrainian pull survives the language-axis projection,
and the cut is itself a property of the OpenAI embedding space. The
weaker version --- ``the English content signal survives debiasing
with a $5\%$--$13\%$ collapse in four of the five conflicts
($38\%$ on the weak Falklands signal); the contested-language
signals collapse
asymmetrically ($96\%$ Hebrew, $66\%$ Ukrainian,
$51\%$ Spanish; $20\%$--$37\%$ for Arabic, Hindi,
Urdu)'' --- is what the four-embedder panel actually licenses.

\section{Reconciling the four measurement
  instruments}\label{sec:four_instruments}

The four-way disagreement set up at the end of
Chapter~\ref{ch:experiments} resolves once each instrument is read
for what it actually measures. Read this way, the disagreement is
itself the methodological contribution.

\emph{Anchor cosine} measures topical similarity between a response
and an own-language Wikipedia article the model has plausibly read in
training. A high English diagonal pull says: when prompted in
English, the response is topically closer to the English Wikipedia
treatment of Crimea than to the Russian or Ukrainian treatments. This
is a statement about \emph{which corpus the response sounds like}, not
about \emph{which side the response takes}. English Wikipedia
treatment of Crimea is a third stance, neither pro-Russian nor
pro-Ukrainian, and a response that sounds like it can be
content-neutral on the contestation itself.

This carries a terminological consequence for the thesis's central
construct. \citet{oeberst2020wikipedia}'s ingroup bias is
resemblance to \emph{the in-group's own account} --- it presumes
the language indexes a conflict party. On the Russo-Ukrainian,
Israel--Palestine, India--Pakistan, and Taiwan families, English
speakers are not a party, so the English diagonal --- the largest
number in the thesis --- measures \emph{own-language corpus pull},
not ingroup bias in Oeberst's party sense; only the Russian,
Ukrainian, Hebrew, Arabic, Hindi, Urdu, and Chinese diagonals carry
the party reading. The one family where the English diagonal
genuinely is an Oeberst-sense ingroup measurement --- the
Falklands, where English is the language of a disputing party ---
is precisely where the pull collapses to $+0.05$. That inversion is
consistent with the corpus-pull reading of the cosine instrument:
what the diagonal responds to is how distinctive the own-language
treatment of the event is, not how invested the language community
is in the dispute, and the Falklands' English-language treatment is
close to the global English mainstream. The thesis keeps
``ingroup pull'' as the operational name for the diagonal, but its
interpretation as in-group \emph{siding} is carried by the stance
instruments, not by cosine.

\emph{Paired-edit stance projection}, by construction, measures the
second question. The axis is built from minimal-edit sentence pairs
that share vocabulary and structure and differ only in which agent
the sentence assigns the action to; projecting a response onto that
axis returns a signed scalar that says which side's framing
vocabulary the response is using, controlled for topic overlap. The
Russo-Ukrainian gap of $+0.033$ says that the
cross-language framing gap is small on this corpus --- smaller than
on India--Pakistan ($+0.131$) or Israel--Palestine
($+0.117$) --- even though the pair's cosine pull is substantial,
mid-ranking among the five conflicts. The stance axis is
finding less framing divergence between Russian and Ukrainian
responses than between Hindi and Urdu responses, or between Hebrew
and Arabic responses, despite Russian and Ukrainian indexing two
internal-narrative communities directly opposed to each other on this
conflict.

\emph{The Sonnet prose-stance judge on imaginary content} measures
something different again. The imaginary-conflict pilot removed the
Wikipedia anchor; the stance judge then coded each response on the
$-2/{+}2$ axis of \S\ref{sec:exp_004} (positive Russian-ward,
negative Ukrainian-ward). Per-language
means EN $-0.092$, RU $+0.010$, UK
$-0.170$ say that, on an event with no anchor for the model
to lean on, the directional signal in free-form prose is small and
tilts Ukrainian-ward overall, not Russian-ward. This is the cleanest
falsification of the literal national-stance hypothesis in the
thesis: when the model cannot defer to a training-corpus treatment of
the event, the residual directional signal does not align with the
own-language nation.

\emph{The discrete-flip MCQ measurement} of
\citet{li2024borderlines} reports yet another direction on the same
Crimea item: vanilla GPT-4 selects the Russian-claimed answer when
asked in Russian and the Ukrainian-claimed answer when asked in
Ukrainian. The MCQ measurement forces a binary commitment that the
prose-generation measurements do not.

The four instruments answer four different questions: which corpus
the response resembles; which side's framing vocabulary it uses; what
direction its prose tilts when no anchor is available; and which
binary commitment it makes when forced to pick. The right reading is
not to pick a winner but to read all four together. On Russo-Ukrainian
content, frontier models reproduce a recognisable English-Wikipedia
\emph{topic} regardless of prompt language, use mostly-shared
\emph{framing vocabulary} across Russian and Ukrainian, tilt slightly
Ukrainian-ward in \emph{unanchored prose}, and flip Russian-ward
under \emph{forced-choice}. The single-instrument national-stance
question is malformed for this kind of model behaviour; the bias is
fully described only by the four-way metric panel.

One rival reading must be bounded before this reconciliation can
carry weight: the four legs differ not only in instrument but in
corpus and model --- two run on the real-event corpus, the judge on
the imaginary corpus, and the MCQ leg is an external result on
vanilla GPT-4, which is not in the matrix and was never run on
these prompts. Part of the disagreement could therefore be
corpus or model rather than instrument. The cross-checks bound how
far that reading goes: the stance projection replicates its
compression across the five-conflict sweep, and the ladder judge
arm reproduces the Ukrainian-ward direction on \emph{real} events
--- though under the same judge model; the second-judge check ran
on the imaginary corpus --- so the instrument-difference
reading holds wherever two instruments were run on the same
corpus. The forced-choice leg is the exception --- it has not been
run on this corpus or these providers, and the four-instrument
panel treats it as a literature anchor whose replication here is
future work, not as a measurement of this corpus.

A scope condition follows directly. The MCQ behaviour
\citet{li2024borderlines} measures requires a fixed answer space and
a forced commitment; the cosine behaviour requires a Wikipedia
anchor; the stance-projection behaviour requires a paired-edit
lexicon. Each instrument has prerequisites that determine where it
applies. Reporting a single ``the model is biased toward
$X$'' headline conceals this.

\section{Two modes of bias on imaginary
  content}\label{sec:two_modes}

Removing the anchor sharpens a second observation. On the imaginary
Russo-Ukrainian conflict, two questions in the bank dissociate
content from direction in opposite ways.

The free-form prose question (``write a 150-word news report'')
produces the highest cross-lingual content divergence in the bank
($2.86\times$ same-language baseline) but the lowest stance
direction (mean $0.00$). Models confabulate different
incidents in different languages without picking a side: the same
fictional bridge is placed over three different rivers in three
languages by Qwen-Plus, with Russian-state phrasing appearing in the
Russian version only. The direct-attribution question (``who is
responsible?'') has the inverse: low content divergence
($1.73\times$) and the strongest directional pull (UK mean
$-0.56$). Models hedge with the same shape across all three
languages but lean Ukrainian-ward when they commit.

Bias on imaginary content therefore has two distinct shapes ---
\emph{content divergence without direction} in free-form prose and
\emph{directional commitment without content divergence} in
attribution --- and the literal national-stance hypothesis collapses
both into one. The same model can be biased on the same fictional
event in opposite ways depending on the question form. The
Cohere \verb|command-r7b-arabic-02-2025| swing --- $+0.47$
(Russian-ward) under Russian prompts to $-0.60$ (Ukrainian-ward)
under Ukrainian prompts, a $1.07$-point cross-language
swing on a single fictional incident --- shows the cell-level
direction can be large, and in this one provider it does align with
the own-language national framing. What the matrix as a whole shows
is that no such alignment survives averaging: the per-language
means are near zero and tilt slightly Ukrainian-ward under
\emph{all three} prompt languages. The literal national-stance
hypothesis fails as a description of the provider matrix while
remaining true of an isolated 7B provider --- a scope distinction
the single mean conceals.

\section{Candidate mechanisms for the cross-lingual
  pull}\label{sec:mechanisms}

The cross-lingual ingroup pull is consistent with at least three
non-exclusive mechanisms; the experiments in this thesis cannot
discriminate among them.

\emph{Training-corpus density} is one. The English-language web is
overrepresented in the training corpora of every provider in the
matrix, and the English Wikipedia article on Crimea is among the
best-indexed treatments of the topic, giving an English prompt the
densest possible target to reproduce.
\citet{zhao2025makieval} have argued that multilingual LLM behaviour
shifts toward English norms even when prompted in other languages ---
an ``English Activation'' effect --- and this offers a mechanistic
story for the cosine pull that does not require the model to hold
ideological preferences at all. The four-embedder panel weakens this
reading slightly --- the OpenAI-specific magnitude ordering suggests
part of the EN ${\gg}$ everything else gap is an embedder property,
not a generator property --- but does not refute it.

\emph{Cross-lingual transfer between Russian and Ukrainian} is a
second. The two languages share Cyrillic script, share a substantial
fraction of subword tokens (\citet{qi2023crosslingual} reports
$0.76$ overlap for BLOOM's tokenizer; the embedder tokenizers
measured in \S\ref{sec:exp_017} sit at $0.31$ to $0.53$), and the
larger Russian web corpus likely supplies a substantial fraction of
the training signal for Ukrainian generation. The unsupervised
clustering finding that the Russian and Ukrainian anchors and both
languages' responses fuse into a single joint cluster on Crimea,
Maidan, and Bandera (\S\ref{sec:exp_016_anchor_projection}) is
consistent with this mechanism, though as a symmetric fusion it
cannot by itself establish the direction of transfer; the
directional evidence is the asymmetric debiasing collapse of
\S\ref{sec:exp_001_debiased}. Critically, the tokenizer-overlap
control rules out subword-vocabulary leakage as the dominant
explanation: at matched cluster resolution under the
lowest-overlap embedder, which questions fuse changes but the total
fusion magnitude does not. What is left is a content-level
phenomenon, not a tokenizer artefact.

\emph{Pre-training reflection of Wikipedia ingroup bias} is the
third. \citet{oeberst2020wikipedia} establish that Wikipedia
articles on intergroup conflicts show systematic ingroup-favourable
framing in the language of the editing community. A model that
reproduces the distribution of its training text on contested events
will reproduce this framing, with no further preference required.
The Falklands near-null cosine result is consistent with this
mechanism --- the Falklands edition gap in
\citet{oeberst2020wikipedia}'s data is among the weakest in their
sample --- and is the cleanest piece of evidence in this thesis for
the source-side mechanism. \citet{smirnov2026language} document the
same phenomenon at the document level on a single Ukrainian
civil-society text, and the multi-event multi-provider extension in
this thesis can be read as a scaling test of that observation.

These three mechanisms make overlapping predictions on cosine and
diverge on the other instruments. Discriminating among them is
outside the scope of the present work and is flagged in
\S\ref{sec:future_work}.

\section{YandexGPT and structural inheritance of state-aligned
  guardrails}\label{sec:yandex_limit}

The YandexGPT finding sits orthogonal to the cosine story. Across
$270$ Russo-Ukrainian prompts the provider returned the same
Russian content-filter boilerplate, with no response in any cell. The
cosine pipeline reads this as null; the unsupervised clustering
reads it as a tight four-cluster pure-provider pocket, which is the
embedding-space signature of categorical content filtering.

Read alone, this is a categorical anomaly. Read alongside
\citet{urman2025silence}'s finding that Western chatbots refuse
roughly $90\%$ of Russian-language Putin-related queries
against under $20\%$ of English-language equivalents, it is
the limit case of the same mechanism: structural inheritance of
state-aligned content-moderation templates, visible across the
multilingual matrix at varying intensities. The five-conflict
extension reported in \S\ref{sec:exp_011} sharpens the reading:
refusal is steepest where Russian state narratives are most
directly implicated --- total on the imaginary Russo-Ukrainian
event, vestigial on the Taiwan strait, with the interest ordering
read descriptively from the rates rather than fitted to an ex-ante
proxy --- and within events the one-sided-defence slots are always
filtered first. A filter graded by state interest and typed by
question form is not a blunt domain block; it is a policy with a
topology, and the embedding-space signature recovers that topology
without access to the policy itself.

The unsupervised method does work here that cosine cannot: the
graded Yandex refusal pocket, the Mercury-2 Crimea avoidance, and
the small-open-weight stylistic fingerprint are all visible only
because the clustering runs alongside the cosine. The methodology
is therefore best framed as a two-instrument minimum: a magnitude
metric for the headline pull, and a structure-recovery metric for
what the magnitude averages away.

\section{Threats to validity}\label{sec:threats}

\subsection{Pilot-scale samples for the stance-judge
  arm}\label{sec:threat_pilots}

The Sonnet stance-judge code from \S\ref{sec:exp_004_metric1} runs on
the imaginary-conflict pilot at $n{=}3$ per cell. The headline
direction --- Ukrainian-ward overall, neutral on free-form prose,
strongest on q01 attribution --- survives the small sample because
the per-language means separate by several times the within-cell
standard deviation. Subtler effects, in particular the
asymmetric-meta-cognition observation that two providers flag the
fictional event in only certain languages, do not. Promoting the
pilot to $n{=}10$ would be a low-cost extension; it was not run, and
the arm should be read accordingly.

\subsection{Single-judge stance coding}\label{sec:threat_single_judge}

The stance code was generated by Sonnet 4.6 alone; the
pre-specified plan called for a second judge for inter-rater
agreement. The disagreement between the Sonnet stance direction
(Ukrainian-ward) and the cosine direction (English-ward) cannot be
disentangled from a potential Sonnet-side latent preference on
Russo-Ukrainian content without a second judge. The event-ladder
judge arm (\S\ref{sec:exp_007_judge}) does not close this gap: it
reproduces the Ukrainian-ward direction on a second corpus and a
different pair of judged providers, but its judge is again Sonnet
4.6, so every directional judge signal in this thesis passes
through a single judge model --- and that arm carries its own
settled-baseline calibration failure besides. The Mercury
topic-avoidance coding of \S\ref{sec:exp_003} is Sonnet-coded under
the same limitation.

The pre-specified remediation --- coding a $30$-response stratified
sample under a second frontier judge from a different provider and
reporting agreement on the directional sign --- was carried out with
OpenAI GPT-4o under the locked rubric, verbatim, at temperature
zero. Of the $28$ responses both judges scored, directional-sign
agreement is $24/28$ ($86\%$, Cohen's $\kappa = 0.78$) with
\emph{no} sign inversions anywhere in the sample: every
disagreement is GPT-4o neutralising a Sonnet nonzero, in both
directions ($3\times$ negative${\to}0$, $1\times$
positive${\to}0$). Exact five-point agreement is lower ($54\%$)
because GPT-4o codes Russian-ward positives one step more extreme;
stance \emph{intensity} is judge-dependent while stance
\emph{direction} is not. The Ukrainian-ward direction is therefore
not attributable to a Sonnet-side latent preference. What remains
open is narrower: the Mercury coding and the ladder judge arm are
still single-judge, and the settled-pair calibration failure of
\S\ref{sec:exp_007_judge} is untouched by this check.

\subsection{Native-speaker review pending}\label{sec:threat_native_review}

The Russian and Ukrainian translations of the prompt bank and the
Tarasenko Bridge fictional place names (Yarov-Kut, Pridvinsk) have
not been reviewed by native speakers. Qwen-Plus on q02 flagged
Yarov-Kut as mapping to a real settlement in Kursk Oblast; the
finding is documented but the place name has not been remediated.
Native-speaker review of the full prompt bank and the fictional
place names is a prerequisite for any publication-track extension.

\subsection{Frontier-scoped claims}\label{sec:threat_scope}

The claim that provider identity is not a structural axis of the
embedding space is empirically scoped to the twelve providers that
behave homogeneously on the embedding geometry --- an operational
set, not a parameter threshold (\S\ref{sec:exp_016}). The two
locally hosted open-weight 7B models --- ALLaM and TAIDE --- retain
stylistic signatures that the
other providers do not, and their same-language baseline
divergence runs $2$--$4\times$ the frontier rate; the API-served
\verb|command-r7b|, also 7B, sits inside the frontier band on
geometry yet carries the largest stance swing of \S\ref{sec:exp_004},
so homogeneity on one instrument does not transfer to
another. A residual confound should be named: the two fingerprinted
models are served through a local quantized Ollama stack while
every other provider serves full-precision weights over an API, so
what the design isolates is the serving stack, not parameter count
or open-weight status --- aggressive quantization alone can
depress output diversity in exactly this way. ``Providers do not
structurally matter'' should be read as ``the geometrically
homogeneous providers do not structurally matter on the (question
$\times$ language) axis.''

\subsection{Anchor size is confounded with event recency}\label{sec:threat_anchor_length}

The Russia--Ukraine 2022 cell carries both the lowest diagonal pull
in the ladder corpus and, by a wide margin, its longest and most
Russian-skewed Wikipedia anchors, so anchor length and event recency
were not separated in the original design. The pre-specified
remediation --- re-running against anchors truncated to the shortest
in their conflict family, which requires only re-embedding --- was
carried out (\S\ref{sec:exp_007_cosine}). The confound is now
quantified rather than open: roughly a third of the 2022 deficit
was anchor length ($+0.075 \to +0.111$ under matching), and the
residual dip survives as the family minimum. Two caveats remain.
First, family-minimum matching is aggressive where a family
contains a stub article --- the China--Japan anchors truncate to 72
tokens --- so a percentile-floor variant would be a gentler control.
Second, truncation changes anchor \emph{content}, so matched values
are comparable only within the matched run, not to the original
gradient's levels. A residual confound the matching cannot touch:
the 2022 and 2023 anchor leads were substantially written after
several providers' training cutoffs, so a model cannot resemble
anchor content that postdates its knowledge --- a mechanism that
depresses live-event diagonals independently of framing and works
against detecting any true recency gradient. Cutoff dates are not
published for every provider in the matrix, so this is noted
rather than quantified.

\subsection{Machine-translated ladder prompts}\label{sec:threat_mt}

Apart from English, and apart from the Russian and Ukrainian prompts
inherited from the case-study bank, every prompt in the event-ladder
corpus is machine-translated and unreviewed by a native speaker. The
ladder covers sixteen languages, including several --- Greek,
Turkish, Japanese, Czech, Slovak, Norwegian, Swedish --- where the
project has no internal reading capacity at all. A translation that
weakens or strengthens the loaded surface form of a question would
move that language's cell without any model behaviour changing.
The cross-family comparisons in \S\ref{sec:exp_007} should be read
with this in mind: the within-family temporal comparisons are the
better-controlled contrast, because the same translation pipeline
applies across a family's ladder, while between-family magnitude
rankings inherit whatever per-language translation quality varies.
Native review is the hardening gate for this arm.

\subsection{English-fitted stance axes}\label{sec:threat_stance_axis}

The paired-edit stance axes are fitted from English seed sentences
(\S\ref{sec:stance_axis}) and every non-English response is
projected onto them through the embedder's cross-lingual alignment.
If that alignment quality varies across language pairs, part of the
between-family gap differences in \S\ref{sec:exp_021} and
\S\ref{sec:exp_007_stance} reflects alignment rather than framing
behaviour; the max-minus-min gap statistic additionally favours
three-language families over the two-language Taiwan and Falklands
families. The within-family cross-language gaps, computed under a
single axis, are the safer contrast, and the Russo-Ukrainian floor
--- the thesis's main stance-arm claim --- is a within-comparison
result that survives both caveats. A further unrun sensitivity
check applies to cosine and stance aggregates alike: the two
forced-framing questions (q07, q08) enter every aggregate alongside
the spontaneous questions, and excluding them was not tested. The between-family
\emph{ranking} should be read as exploratory until the axes are
re-fitted from native-language seed pairs, which the present design
does not include.

\subsection{Wikipedia anchors are not culturally
  neutral}\label{sec:threat_anchors}

The cosine pipeline treats own-language Wikipedia articles as
anchors. \citet{naous2024beer} argue that even
own-language Wikipedia carries Western-centric framing on
non-Western topics. The English-control design of this thesis is
partially answerable on this point --- a finding that the English
diagonal pull is content-driven survives the projection onto the
language subspace, suggesting the English signal is something more
than language similarity --- but the anchors themselves remain a
limit on how much the cosine pull can be read as a measurement of
neutral cross-lingual divergence.

\section{Implications}\label{sec:implications}

Two operational implications follow.

For \emph{multilingual deployment} the relevant question is not
``which provider is neutral?'' but ``what does neutrality mean
when the prompt language is part of the framing?'' The cosine evidence shows
the prompt language --- not the
provider choice --- is the dominant determinant of which corpus
the response resembles: frontier cross-provider agreement approaches
within-provider self-agreement, and provider identity surfaces only
marginally as a structural axis of the embedding space --- a claim
scoped to frontier-scale providers
(\S\ref{sec:threat_scope}). A deployment that wants to
control the framing of model output on contested topics has to
control the prompt language; switching frontier providers does not.

For \emph{evaluation of multilingual LLM bias} the four-way metric
panel argues against single-instrument headlines. Anchor cosine,
unsupervised clustering, paired-edit stance projection, and discrete
forced-choice each measure a different facet of cross-lingual
behaviour, and a paper that reports only one is incomplete. The
methodology contribution of the thesis is to demonstrate this on a
single corpus and to map the conditions under which each instrument
is informative.

\section{Future work}\label{sec:future_work}

Two of the extensions this section previously listed --- the
length-matched anchor re-run and the second stance judge --- were
carried out and are now reported in
\S\ref{sec:threat_anchor_length} and
\S\ref{sec:threat_single_judge}. Three extensions remain, ordered
by how much they would strengthen the existing claims.

\emph{Native review of the ladder prompts.} The sixteen-language
event-ladder corpus (\S\ref{sec:threat_mt}) is machine-translated
throughout. Native review of even a subset --- the highest-magnitude
families, Greek--Turkish and Hindi--Urdu --- would convert the
between-family stance ranking from exploratory to load-bearing.

\emph{Settled-pair judge calibration and full-corpus second
coding.} The second-judge check establishes directional agreement
on a stratified sample; extending it to the full $630$-response
corpus, and calibrating both judges against the settled pairs where
the ladder judge arm currently miscalibrates
(\S\ref{sec:exp_007_judge}), would close the remaining judge-side
gaps. A percentile-floor variant of the anchor matching belongs in
the same hardening pass
(\S\ref{sec:threat_anchor_length}).

\emph{The four-way metric reconciliation extended to the full
provider matrix.} The discrete-flip MCQ measurement of
\citet{li2024borderlines} on Crimea, Donbas, Sevastopol, and South
Ossetia ran on a small provider set; running the same MCQ items on
the provider matrix used here closes the cosine
vs.\ stance vs.\ judge vs.\ MCQ comparison at single-provider
granularity and either reproduces the binary--prose disagreement at
scale or isolates it to GPT-4.

Three further directions are outside the scope of this thesis but
follow naturally from its design. \emph{Provider-country bias as an
axis orthogonal to language} --- holding the prompt language fixed
while varying the provider's regulatory home --- is the obvious
complement to the language-conditional design used throughout, and
a contested event on which the provider countries themselves are the
disputing parties would separate the two axes cleanly. \emph{A
French-language control on a European provider} would test whether
the English-ward pull reported here is specific to English or is a
general flagship-language effect, which the present matrix cannot
distinguish because its non-English languages are all languages of a
disputing party. \emph{Three-way rather than binary framing spaces},
on disputes where the contest is not adequately described by two
poles, would test whether the paired-edit stance construction of
\S\ref{sec:stance_axis} generalises beyond the bipolar case it assumes.

\section{Conclusion}\label{sec:conclusion}

The four research questions of \S\ref{sec:rqs} can now be answered
directly.

\emph{RQ1 --- ingroup pull.} The answer is per-language.
Decisively yes for English: matrix mean $+0.175$, and still
$[+0.13, +0.18]$ after the language-axis control --- largely
content-driven. Yes in raw terms for Russian ($+0.043$) and
Ukrainian ($+0.026$), whose question-resampled intervals exclude
zero --- though under the stricter anchor-block resampling of
\S\ref{sec:exp_001} the Russian interval touches zero, so the raw
contested-language pulls are resolved only marginally. After the
control, the Russian residual ($+0.031$, interval
$[+0.005, +0.06]$) is barely positive and the Ukrainian residual
($+0.009$, interval $[-0.02, +0.03]$) is indistinguishable from
zero --- two thirds of
the Ukrainian pull was lexical, and the remainder is within
noise. All of this is in the OpenAI embedding space; other
embedders draw the lexical-content line elsewhere.

\emph{RQ2 --- generality.} The sign of the pull survives across the
provider matrix and across conflict families: five contested
families in the scale-up, eight bilateral disputes and 31 events
across a millennium in the ladder corpus, with two settled-dispute
baselines at the floor of the cosine arm in raw pull, below every
contested family under length-matched anchors --- though at the
2022 cell that separation is a point estimate of $+0.007$, and the
settled floor is partly mechanical, since the settled pairs'
mutually intelligible languages compress their attainable pull
(\S\ref{sec:exp_007}).
The exceptions are informative rather than damaging: the
Taiwan-strait Chinese cell is marginally negative, the 7B
open-weight providers carry stylistic signatures that scope
provider-agnostic claims to frontier models, and the Russian
provider refuses categorically --- itself a regime-level finding.
The recency gradient documented for human editors does not
transfer: bias direction is inherited without its temporal
structure.

\emph{RQ3 --- robustness.} The sign is not an embedder artefact: it
holds under four embedders from four ecosystems, and the
Russian-trained embedder does not inflate the Russian diagonal. The
magnitude ordering EN ${\gg}$ RU ${>}$ UK, by contrast, is
OpenAI-specific and is treated as an embedder property throughout.
Subword-tokenizer overlap is ruled out as the dominant cause of the
Russian--Ukrainian cluster fusion: fusion redistributes but does
not shrink under the lowest-overlap tokenizer.

\emph{RQ4 --- what cosine misses.} A great deal, and systematically.
The stance-axis projection demotes the Russo-Ukrainian pair from
mid-ranking on cosine to the floor of the contested families; on
imaginary content the bias splits into content divergence without
direction (free-form prose) and directional commitment without
content divergence (attribution); and refusal and topic avoidance
are categorical behaviours the cosine metric cannot see at all. The
four instruments disagree because they measure four different
things, and the single-instrument question ``which side is the
model biased toward?'' is malformed for this class of model
behaviour.

The overall conclusion is twofold. Substantively, of the two
variables a deployer controls, the prompt language --- not the
provider --- is the dominant determinant of what a contested-event
response resembles at frontier scale (question identity,
which the user supplies, carries more cluster structure than
either), and what the language
controls is which corpus the response resembles rather than which
side the response takes. Methodologically, no single instrument
suffices to characterise cross-lingual bias in multilingual LLMs;
the minimum honest report is a panel, and this thesis supplies a
worked example of what such a panel looks like and where each of
its instruments stops being informative.
